\documentclass[12pt]{ximera}
\title{Homework 6: Fair Division --- Divider-Chooser and Lone-Divider}
\begin{document}
\begin{abstract}
Sections 3.1--3.3: continuous versus discrete fair division, computing a fair share
from one player's value system, the lone-divider method (including a standoff that
forces a re-cut), and divider-chooser applied to a half-veggie, half-meatball
sandwich.
\end{abstract}
\maketitle

\section*{Sections 3.1--3.3: Fair Division Games}

\begin{problem}
The four Dalton brothers want to divide the loot after robbing a bank vault. Their
names are D1, D2, D3, and D4. For each of the following scenarios, state whether the
assets are continuous, discrete, or a combination of both.

\begin{prompt}
\textbf{(a)} The vault contained a Grant Wood painting of a couple wielding a pitchfork
and the map to Scarface's treasure chest.

\begin{multipleChoice}
\choice{Continuous}
\choice[correct]{Discrete}
\choice{A combination of both}
\end{multipleChoice}

\begin{feedback}
Discrete. Both items are single, indivisible objects --- you cannot give one brother
$40\%$ of a painting without destroying it.
\end{feedback}
\end{prompt}

\begin{prompt}
\textbf{(b)} The vault contained an authentic sheriff's badge and $3$ barrels of
gunpowder.

\begin{multipleChoice}
\choice{Continuous}
\choice{Discrete}
\choice[correct]{A combination of both}
\end{multipleChoice}

\begin{feedback}
A combination. The badge is a single indivisible object (discrete), while gunpowder
is a bulk substance that can be split into any quantity you like (continuous).
\end{feedback}
\end{prompt}

\begin{prompt}
\textbf{(c)} The vault contained \$205 in solid gold, a dozen pounds of chewing
tobacco, and a barrel of Bourbon.

\begin{multipleChoice}
\choice[correct]{Continuous}
\choice{Discrete}
\choice{A combination of both}
\end{multipleChoice}

\begin{feedback}
Continuous. All three --- gold, tobacco, and Bourbon --- are bulk goods that can be
divided into arbitrarily sized portions without ruining them.
\end{feedback}
\end{prompt}
\end{problem}

\begin{problem}
Suppose the vault does contain \$205 in solid gold, a dozen pounds of chewing tobacco,
and a barrel of Bourbon. If D1 values the tobacco at \$195 and the Bourbon at \$160,
then what is the minimum value of a fair share for him?

$\$\answer{140}$

\begin{hint}
First find what D1 thinks the \emph{whole} vault is worth, then remember there are
four brothers splitting it.
\end{hint}

\begin{feedback}
In D1's value system the vault is worth
\[ \$205 + \$195 + \$160 = \$560. \]
With four brothers dividing it, a fair share is at least one quarter of that:
\[ \frac{\$560}{4} = \$140. \]

Note that a fair share is defined \emph{relative to the player's own value system}.
The other brothers will each compute a different dollar figure, and all of them are
correct for that brother.
\end{feedback}
\end{problem}

\begin{problem}
Suppose Dora has packed a picnic lunch to share with Boots, Tico, and Swiper (call
them D, B, T, and S). They decide to use the lone-divider method.

\begin{prompt}
\textbf{(a)} Suppose their value systems look like this:

\begin{center}
\begin{tabular}{r|c|c|c|c|}
& $s_1$ & $s_2$ & $s_3$ & $s_4$ \\\hline
Dora   & 35\% & 20\% & 25\% & 20\% \\\hline
Boots  & 25\% & 25\% & 25\% & 25\% \\\hline
Tico   & 10\% & 34\% & 40\% & 16\% \\\hline
Swiper & 20\% & 27\% & 48\% &  5\% \\\hline
\end{tabular}
\end{center}

Which of the characters was selected to be the divider?

\begin{multipleChoice}
\choice{Dora}
\choice[correct]{Boots}
\choice{Tico}
\choice{Swiper}
\end{multipleChoice}

\begin{hint}
The divider cuts the lunch into pieces that are all of equal value \emph{to them}.
Whose row gives that away?
\end{hint}

\begin{feedback}
Boots. The divider's job is to cut the lunch into $N$ pieces they consider exactly
equal, so the divider's row must read $25\%$ across the board --- which is Boots's row.

Everyone else is a chooser, and their rows reflect what they happen to think of the
pieces Boots produced.
\end{feedback}
\end{prompt}

\begin{prompt}
\textbf{(b)} Determine two different fair divisions of the picnic lunch.

First, the bid lists. A chooser bids on every share worth at least $25\%$ to them.

Dora bids on $s_1$ and $s_{\answer{3}}$. \quad
Tico bids on $s_2$ and $s_{\answer{3}}$. \quad
Swiper bids on $s_2$ and $s_{\answer{3}}$.

\smallskip
In both fair divisions, Dora must receive $s_{\answer{1}}$ and Boots must receive
$s_{\answer{4}}$.

\begin{hint}
Tico and Swiper bid on exactly the same two shares, so between them they must take
those two. What does that leave for Dora?
\end{hint}

\begin{feedback}
The bid lists are
\[
\text{Dora}: \set{s_1, s_3}, \qquad
\text{Tico}: \set{s_2, s_3}, \qquad
\text{Swiper}: \set{s_2, s_3},
\]
since those are the shares each values at $25\%$ or more.

Tico and Swiper bid on the same pair $\set{s_2, s_3}$, so between them they must take
$s_2$ and $s_3$ --- there is no other option. That leaves Dora with $s_1$, which is
fine since $s_1$ is on her bid list (she values it at $35\%$). Boots, as the divider,
takes whatever remains: $s_4$, which he values at $25\%$.

The two fair divisions are:
\begin{itemize}
\item Dora $s_1$, Tico $s_2$, Swiper $s_3$, Boots $s_4$;
\item Dora $s_1$, Tico $s_3$, Swiper $s_2$, Boots $s_4$.
\end{itemize}
Every player receives a share they value at $25\%$ or more, so both are fair.
\end{feedback}
\end{prompt}

\begin{prompt}
\textbf{(c)} Suppose instead their value systems look like this:

\begin{center}
\begin{tabular}{r|c|c|c|c|}
& $s_1$ & $s_2$ & $s_3$ & $s_4$ \\\hline
Dora   & 15\% & 20\% & 45\% & 20\% \\\hline
Boots  & 25\% & 25\% & 25\% & 25\% \\\hline
Tico   & 10\% & 34\% & 40\% & 16\% \\\hline
Swiper & 20\% & 21\% & 54\% &  5\% \\\hline
\end{tabular}
\end{center}

Can you distribute all four of these shares without cutting them?

\begin{multipleChoice}
\choice{Yes}
\choice[correct]{No --- two choosers bid on only the same single share}
\end{multipleChoice}

\begin{feedback}
No. The bid lists are now
\[
\text{Dora}: \set{s_3}, \qquad \text{Tico}: \set{s_2, s_3}, \qquad
\text{Swiper}: \set{s_3}.
\]
Dora and Swiper each consider \emph{only} $s_3$ to be worth a fair share, and there is
just one $s_3$. No assignment of whole shares can satisfy both, so we have a standoff.
\end{feedback}
\end{prompt}

\begin{prompt}
\textbf{(d)} Describe a possible fair division by cutting some of the shares.

\begin{freeResponse}
\end{freeResponse}

\begin{feedback}
The standard resolution: settle the players who are not part of the standoff first,
then re-divide the rest between the two who are.

\begin{enumerate}
\item Give $s_2$ to Tico. He values it at $34\% \ge 25\%$, so this is a fair share for
      him, and he is the only chooser outside the standoff.
\item Give $s_1$ to Boots. As the divider he values every share at $25\%$, so any
      leftover share is fair for him.
\item Recombine the remaining pieces --- $s_3$ and $s_4$ --- into one lot, and have
      Dora and Swiper divide it between themselves using the divider-chooser method.
\end{enumerate}

Check that step 3 leaves both of them satisfied. Dora values $s_3 + s_4$ at
$45+20 = 65\%$ of the whole lunch, and Swiper values it at $54+5 = 59\%$. Each needs
only $25\%$ to be satisfied, and divider-chooser guarantees each at least half of what
they think the lot is worth --- at least $32.5\%$ for Dora and at least $29.5\%$ for
Swiper. Both comfortably clear $25\%$, so the division is fair for all four players.
\end{feedback}
\end{prompt}
\end{problem}

\begin{problem}
Adam and Joel are sharing the following sandwich.

\begin{center}
\includegraphics[width=0.7\textwidth]{sandwich.jpg}
\end{center}

This sandwich is half vegetarian (call it V) and half meatball (call it M). Adam likes
the veggie part only half as much as the meatball part, but Joel likes both equally.
They've agreed that Adam will make \emph{a single cut} to produce two fair shares in
his opinion, then Joel will choose.

\begin{prompt}
\textbf{(a)} What vertical cut can Adam choose to produce fair shares for himself?

Adam should cut through the meatball half, leaving one share consisting of all of V
together with $\answer{1/4}$ of M, and the other share consisting of the remaining
$\answer{3/4}$ of M.

\begin{hint}
Assign values: say the whole M half is worth $2$ units to Adam and the whole V half is
worth $1$ unit. What is his total, and what does each fair share have to be worth?
Then check whether a cut inside the V half could ever work.
\end{hint}

\begin{feedback}
Set up Adam's value system. Let the M half be worth $2$ units and the V half $1$ unit,
since he likes veggie half as much. His total is $3$ units, so each share must be
worth $1.5$ units.

First check whether the cut could be inside the V half. That would give one share
consisting of part of V only --- worth at most $1$ unit --- which is less than $1.5$.
So the cut must fall inside the M half.

Let the first share be all of V plus a fraction $y$ of M. Adam values it at $1 + 2y$,
and the leftover at $2(1-y)$. Setting them equal:
\[ 1+2y = 2-2y \quad\Longrightarrow\quad 4y = 1 \quad\Longrightarrow\quad y = \tfrac14. \]

So Adam cuts one quarter of the way into the meatball half:
\begin{itemize}
\item Share 1 $=$ all of V $+$ $\frac14$ of M, worth $1 + 0.5 = 1.5$ units;
\item Share 2 $=$ $\frac34$ of M, worth $1.5$ units.
\end{itemize}
Equal in Adam's eyes. \checkmark
\end{feedback}
\end{prompt}

\begin{prompt}
\textbf{(b)} If the sandwich costs \$12 total, how much is each of the shares Adam has
produced worth to Joel?

Share 1 (all of V $+\frac14$ of M): $\$\answer[tolerance=0.005]{7.50}$ \qquad
Share 2 ($\frac34$ of M): $\$\answer[tolerance=0.005]{4.50}$

\begin{hint}
Joel likes both halves equally, so for him value is just proportional to how much
sandwich a share contains.
\end{hint}

\begin{feedback}
Joel values the sandwich purely by size. Measuring shares as fractions of the whole
sandwich (each half being $\frac12$):
\[
\text{Share 1} = \tfrac12 + \tfrac14\cdot\tfrac12 = \tfrac12+\tfrac18 = \tfrac58,
\qquad
\text{Share 2} = \tfrac34\cdot\tfrac12 = \tfrac38.
\]
At \$12 for the whole sandwich,
\[
\tfrac58 \cdot 12 = \$7.50, \qquad \tfrac38\cdot 12 = \$4.50.
\]
\end{feedback}
\end{prompt}

\begin{prompt}
\textbf{(c)} Which share will Joel choose?

\begin{multipleChoice}
\choice[correct]{Share 1 --- all of V plus a quarter of M, worth \$7.50 to him}
\choice{Share 2 --- three quarters of M, worth \$4.50 to him}
\end{multipleChoice}

\begin{feedback}
Joel takes Share 1, worth \$7.50 to him --- well above the \$6.00 he would consider a
fair half.

This is the divider-chooser method working as intended. Adam is satisfied because both
shares are worth exactly half to him, and Joel does \emph{better} than a fair share
because his values differ from Adam's. The chooser's advantage grows the more the two
players disagree about what things are worth.
\end{feedback}
\end{prompt}
\end{problem}

\begin{problem}
Suppose Adam and Joel \emph{and} their friend Jake hope to share the same sandwich.
Jake only likes the veggie part, so doesn't value the meatball part at all. They've
agreed that Joel will make \emph{exactly two cuts} to produce what he sees as three
fair shares, then Adam and Jake will bid on the shares.

\begin{prompt}
\textbf{(a)} If Joel doesn't know either of the others' preferences, what are the three
shares he creates?

Joel values the sandwich by size alone, so he cuts it into three equal thirds:
\begin{itemize}
\item $s_1 = \answer{2/3}$ of V
\item $s_2 = \frac13$ of V $+$ $\frac13$ of M
\item $s_3 = \answer{2/3}$ of M
\end{itemize}

\begin{hint}
Each third of the sandwich is $\frac13$ of the whole, while each of V and M is
$\frac12$ of the whole. What fraction of the V half does the first third take up?
\end{hint}

\begin{feedback}
Joel values both halves equally, so three fair shares for him are simply three equal
thirds by size.

The V half occupies the first $\frac12$ of the sandwich. The first third ends at
$\frac13$, which is inside V, and $\frac13 \div \frac12 = \frac23$ --- so
$s_1$ is $\frac23$ of V.

The second third runs from $\frac13$ to $\frac23$, crossing the midpoint: it takes the
remaining $\frac13$ of V and the first $\frac13$ of M.

The last third is the remaining $\frac23$ of M.
\end{feedback}
\end{prompt}

\begin{prompt}
\textbf{(b)} What is the value of each share from each friend's perspective? Give each
as a percentage of that person's total, rounded to one decimal place.

Adam: $s_1 = \answer[tolerance=0.05]{22.2}\%$, \quad
$s_2 = \answer[tolerance=0.05]{33.3}\%$, \quad
$s_3 = \answer[tolerance=0.05]{44.4}\%$

\smallskip
Jake: $s_1 = \answer[tolerance=0.05]{66.7}\%$, \quad
$s_2 = \answer[tolerance=0.05]{33.3}\%$, \quad
$s_3 = \answer[tolerance=0.05]{0}\%$

\smallskip
Joel: $s_1 = \answer[tolerance=0.05]{33.3}\%$, \quad
$s_2 = \answer[tolerance=0.05]{33.3}\%$, \quad
$s_3 = \answer[tolerance=0.05]{33.3}\%$

\begin{feedback}
\textbf{Adam} values V at $\frac13$ of his total and M at $\frac23$ (the $1:2$ ratio
from Problem 4). So
\[
s_1 = \tfrac23\cdot\tfrac13 = \tfrac29 \approx 22.2\%,
\]
\[
s_2 = \tfrac13\cdot\tfrac13 + \tfrac13\cdot\tfrac23 = \tfrac19+\tfrac29 = \tfrac13 \approx 33.3\%,
\]
\[
s_3 = \tfrac23\cdot\tfrac23 = \tfrac49 \approx 44.4\%.
\]

\textbf{Jake} values only V, so the V half is $100\%$ of his total:
\[ s_1 = \tfrac23 \approx 66.7\%, \qquad s_2 = \tfrac13 \approx 33.3\%, \qquad s_3 = 0\%. \]

\textbf{Joel} values by size, so each third is $33.3\%$.

Each row totals $100\%$. \checkmark
\end{feedback}
\end{prompt}

\begin{prompt}
\textbf{(c)} Determine \emph{three} fair divisions of this sandwich among the three
friends, without any further cutting.

\begin{freeResponse}
\end{freeResponse}

\begin{hint}
A fair share means at least $\frac13 \approx 33.3\%$ in that person's own values.
Write down each person's bid list first --- which shares clear that bar for them?
\end{hint}

\begin{feedback}
Each player needs a share worth at least $33.3\%$ to them, giving the bid lists
\[
\text{Adam}: \set{s_2, s_3}, \qquad \text{Jake}: \set{s_1, s_2}, \qquad
\text{Joel}: \set{s_1, s_2, s_3}.
\]
($s_1$ is only $22.2\%$ to Adam, and $s_3$ is worthless to Jake.)

There are exactly three assignments giving everyone a share from their own list:
\begin{enumerate}
\item Adam $s_3$ ($44.4\%$), Jake $s_1$ ($66.7\%$), Joel $s_2$ ($33.3\%$);
\item Adam $s_2$ ($33.3\%$), Jake $s_1$ ($66.7\%$), Joel $s_3$ ($33.3\%$);
\item Adam $s_3$ ($44.4\%$), Jake $s_2$ ($33.3\%$), Joel $s_1$ ($33.3\%$).
\end{enumerate}

\textbf{Which is most equitable?} Division 3 gives perceived values of $44.4\%$,
$33.3\%$, and $33.3\%$ --- the most nearly equal spread, so it is the most equitable in
the sense of treating the players alike.

Division 1 is the most \emph{efficient}: its perceived values total
$44.4+66.7+33.3 = 144.4\%$, the largest of the three, because it hands each share to
the person who wants it most. It is arguably the better outcome overall, even though
Jake ends up far happier than the others --- which is a reasonable point to argue
either way.
\end{feedback}
\end{prompt}
\end{problem}

\end{document}
